% Propietario conservado: /Users/ruben/Documents/ChatGPT/jueces y controles/REVISION_ARGUMENTAL_APERTURAS_CIERRES_20260915/delivery/ARTICLE_V_ES_ARGUMENTAL_20260915/manuscrito/sections/30_fock_gibbs.tex
% Recuperación documental para el artículo V; RESULTADO_RECUPERADO.
% Propietario: /Users/ruben/Documents/New project/output/REVISION_EDITORIAL_HMT_MD_20260905/01_FUENTES/INTEGRAL/tree/New project/output/ENSAMBLADO_CAUSAL_RESTAURADO_HMT_MD_20260905/fuente/manuscrito/sucesor_102/deltas_masas/realizacion_equivariante_especies_operadores_masicos_publico.tex
% Líneas de origen: 856--932.
% REV02: transporte functorial; CAR/CCR se construyen en 20.
% Correspondencia editorial: gestion/CONSOLIDACION_PAULI_FOCK_REV02.json.
\section{Transporte de las completaciones de Fock por refinamiento}
\label{v:v:rec:fock_gibbs:section:01}

Las completaciones y sus álgebras de operadores están construidas en la
sección~\ref{v:v:sec:completaciones-algebraicas}. El resultado de este tramo
es su compatibilidad con los morfismos de refinamiento: el transporte de
la fibra de una partícula se prolonga a todos los sectores de ocupación.

\subsection{Segunda cuantización de los morfismos de refinamiento}
\label{v:v:rec:fock_gibbs:section:02}

Sea \(\mathsf{Hol}_{\pm}^{\leq1}\) la categoría monoidal de fibras de una
partícula, con morfismos de refinamiento contractivos que conservan la
paridad de holonomía \(\epsilon\in\{+1,-1\}\). Incluye, en particular,
los entrelazadores unitarios. La operación monoidal sobre las fibras es
la suma directa; en las completaciones induce el producto tensorial,
graduado en el sector exterior. La segunda cuantización utiliza las
completaciones de la sección~\ref{v:v:rec:completaciones:section:03}:
\[
 \mathfrak F(H,\epsilon)=
 \begin{cases}
 \mathcal F_+(H),
     &\epsilon=+1,\\[2mm]
 \mathcal F_-(H),
     &\epsilon=-1,
 \end{cases}
\]
y, para un entrelazador contractivo \(T:H\to H'\),
\[
 \mathfrak F(T)=
 \bigoplus_{n\ge0}\operatorname{Sym}^n(T)
 \quad\text{o}\quad
 \bigoplus_{n\ge0}\bigwedge^n\!T.
\]

\subsection{Functorialidad, operadores canónicos y dominios de transporte}
\label{v:v:rec:fock_gibbs:section:03}

Sea \(\mathcal D_\epsilon(H)\) la suma directa algebraica de los sectores
de partículas finitas. En las identidades siguientes, \(a_H^\dagger(f)\)
y \(a_H(f)\) designan los operadores de creación y aniquilación de la
completación elegida: los operadores \(a\) del sector exterior o los
operadores \(b\) del sector simétrico de la
sección~\ref{v:v:pauli:car-construccion}.

\begin{theorem}[Transporte interno de la estadística]
\label{v:v:rec:fock_gibbs:statement:01}
La completación anterior es un funtor de operadores acotados sobre
\(\mathsf{Hol}_{\pm}^{\leq1}\), compatible con refinamiento y
monoidal graduada. En la rama \(\epsilon=-1\) los operadores de creación y
aniquilación satisfacen CAR y los operadores de ocupación cumplen
\(N_s^2=N_s\); en la rama \(\epsilon=+1\) satisfacen CCR. Por tanto, la
exclusión de Pauli y las estadísticas fermiónica y bosónica descienden de la
paridad de holonomía una vez fijada la representación de una partícula.
Para todo entrelazador contractivo \(T:H\to H'\), \(f\in H\) y
\(g\in H'\), se cumplen en \(\mathcal D_\epsilon(H)\)
\begin{align}
 \mathfrak F(T)a_H^\dagger(f)
 &=a_{H'}^\dagger(Tf)\mathfrak F(T),
 \label{v:v:fock:eq:transporte-creacion}\\
 a_{H'}(g)\mathfrak F(T)
 &=\mathfrak F(T)a_H(T^*g).
 \label{v:v:fock:eq:transporte-aniquilacion}
\end{align}
Si \(T\) es isométrico, la segunda identidad, con \(g=Tf\), transporta
también la aniquilación de \(f\). Para \(\|f\|=1\), los operadores de
ocupación satisfacen entonces
\begin{equation}
 N_{H',Tf}\mathfrak F(T)=\mathfrak F(T)N_{H,f},
 \qquad N_{H,f}=a_H^\dagger(f)a_H(f).
 \label{v:v:fock:eq:transporte-ocupacion}
\end{equation}
\end{theorem}

\begin{proof}
Las potencias simétricas y exteriores preservan composiciones e identidades;
en cada grado sus normas están acotadas por \(\|T\|^n\). La contractividad
garantiza así que la suma directa define un operador acotado. La
descomposición de las potencias de una suma directa por sus grados en cada
factor da la compatibilidad monoidal, con el signo de intercambio graduado
en las potencias exteriores.

Las relaciones CAR y la idempotencia son el
teorema~\ref{v:v:rec:completaciones:statement:01}; las relaciones CCR son
\eqref{v:v:pauli:eq:ccr}. Para comprobar su transporte, basta actuar sobre
productos simétricos o exteriores de vectores. Aplicar \(T\) a cada factor
conmuta con insertar el factor \(f\), que se transforma en \(Tf\);
esto prueba \eqref{v:v:fock:eq:transporte-creacion}. La aniquilación suma las
contracciones de cada factor y utiliza
\(\langle g,Tf\rangle=\langle T^*g,f\rangle\); los signos exteriores y
las normalizaciones simétricas coinciden a ambos lados. Se obtiene
\eqref{v:v:fock:eq:transporte-aniquilacion}. Si \(T^*T=I\), poner \(g=Tf\)
y componer las dos identidades demuestra
\eqref{v:v:fock:eq:transporte-ocupacion}. Todas estas operaciones preservan
el dominio de partículas finitas.
\end{proof}

Para morfismos no contractivos, la segunda cuantización se define primero
en \(\mathcal D_\epsilon(H)\); su extensión acotada depende de la cota
uniforme de las normas de todos los grados. Tales morfismos quedan fuera
de \(\mathsf{Hol}_{\pm}^{\leq1}\). En la completación simétrica, si
\(\|T\|>1\), las normas \(\|\operatorname{Sym}^nT\|=\|T\|^n\)
crecen sin cota; en la completación exterior de una fibra finita sólo
hay un número finito de grados no nulos. La identidad general de
aniquilación conserva \(T^*g\); su especialización al mismo vector
transportado corresponde a la condición isométrica indicada.

La distinción se verifica ya en una fibra unidimensional. Si
\(T=tI\), con \(0<t<1\), y \(e\) es unitario, entonces, sobre el estado
de una partícula \(e\), en cualquiera de las dos completaciones,
\begin{equation}
 a(te)\mathfrak F(T)e=t^2\Omega,
 \qquad \mathfrak F(T)a(e)e=\Omega,
 \label{v:v:fock:eq:control-no-isometrico}
\end{equation}
donde \(\Omega\) es el vacío. La identidad general
\eqref{v:v:fock:eq:transporte-aniquilacion} da correctamente
\(\mathfrak F(T)a(T^*te)e=t^2\Omega\). El factor \(t^2\) identifica
exactamente la hipótesis que permite especializar el transporte de la
aniquilación y, después, el de la ocupación.

El resultado establece el transporte algebraico interno. Su identificación
con el teorema relativista de espín--estadística conserva las condiciones
adicionales de representación lorentziana y localidad explicitadas en la
sección~\ref{v:v:rec:completaciones:section:05}.

\clearpage
\section{Realización termodinámica mediante estados de Gibbs y límite de presión}
\label{v:v:rec:fock_gibbs:section:04}

\subsection{Hamiltoniano finito, partición, estado de Gibbs y límite cofinal de la presión}
\label{v:v:rec:fock_gibbs:section:05}

Para un Hamiltoniano comprimido \(H_{X,V}\) sobre un volumen o refinamiento
finito \(V\), el estado térmico está completamente definido por
\[
 Z_{X,V}(\beta)=\operatorname{Tr}e^{-\beta H_{X,V}},
 \qquad
 \rho_{X,V,\beta}=Z_{X,V}(\beta)^{-1}e^{-\beta H_{X,V}}.
\]
El límite termodinámico existe cuando la presión
\[
 p_X(\beta)=\lim_{V\nearrow\infty}
 \frac1{\beta |V|}\log Z_{X,V}(\beta)
\]
existe y es independiente de la sucesión cofinal.

\subsection{Criterio de condensación por saturación de la densidad excitada}
\label{v:v:rec:fock_gibbs:section:06}

La condensación bosónica se
decide por la saturación de la densidad de estados excitados; no se proclama a
partir de la mera existencia de una completación simétrica. Así, el problema
queda expresado mediante una condición matemática verificable para el
Hamiltoniano y el límite, en lugar de una atribución nominal.
